\cleardoublepage

\newrefsection

\chapter{外文翻译}

% 原文：Yao et al., "ReAct: Synergizing Reasoning and Acting in Language Models", ICLR 2023

\section{引言}

人类智能的一个独特特征，是能够将面向任务的行动与语言推理（即内心独白）无缝结合。有理论认为，这种能力在人类认知中扮演着关键角色，能够实现自我调节与策略制定，并维持工作记忆。以厨房做菜为例，在两个具体的操作步骤之间，我们会通过语言进行推理：追踪操作进度（"食材都切好了，该烧水煮锅了"）、处理突发状况或根据实际情况调整计划（"没有盐了，那就用酱油和胡椒粉替代吧"），并判断何时需要获取外部信息（"怎么揉面团？我去网上查一下"）。我们也会通过行动（翻开食谱看做法、打开冰箱检查食材）为推理提供支撑，解答当下的问题（"现在能用现有食材做什么菜？"）。这种"行动"与"推理"的深度协同，让人类即便在从未遇到的场景中，或是面对信息不确定的情况时，也能快速学习新任务，做出稳健的决策与推理。

近期的研究成果表明，在自主系统中融合语言推理与交互式决策具备可行性。一方面，经过合理提示的大语言模型（LLMs）展现出了一种涌现能力 —— 能够生成多步推理轨迹，从而在算术、常识推理和符号推理任务中推导出答案。但这种"思维链"推理是一个静态的黑箱：模型仅依靠自身的内部表征生成推理过程，并未与外部世界建立关联，这使其无法进行动态的反应式推理，也无法更新自身知识，进而容易出现事实幻觉和推理过程中的错误传播问题。另一方面，已有研究探索将预训练语言模型应用于交互式环境中的规划与行动决策，这类研究的核心是通过语言先验预测行动。这些方法通常将多模态观测信息转换为文本，利用语言模型生成特定领域的行动或计划，再通过控制器选择并执行。但它们并未让语言模型对高层目标进行抽象推理，也未借助工作记忆为行动提供支撑。除了这类与简单积木交互的具身任务外，目前尚无研究探索如何将推理与行动以协同的方式融合，应用于通用任务求解；也未验证这种融合相比单独的推理或行动，是否能带来系统性的性能提升。

本研究提出了 ReAct（Reason + Act）范式，这是一种将大语言模型的推理与行动能力相结合的通用范式，可用于解决各类语言推理与决策任务。ReAct 通过提示大语言模型，以交替的方式生成与任务相关的语言推理轨迹和行动指令，让模型能够进行动态推理，从而制定、维护并调整面向行动的高层计划（为行动而推理）；同时，模型还能与外部环境（如维基百科）交互，将外部信息融入推理过程（为推理而行动）。

我们在四类不同的基准数据集上，对 ReAct 和当前最优的基线方法进行了实证评估：问答任务、事实验证任务、文本基游戏任务和网页导航任务。在 HotPotQA 和 FEVER 任务中，模型可调用维基百科 API 进行交互，实验结果显示，ReAct 的性能优于纯行动生成模型，且与思维链推理（CoT）方法的性能相当。而整体性能最优的方法，是将 ReAct 与 CoT 相结合的融合策略 —— 该策略能让模型在推理过程中同时利用内部知识和外部获取的信息。在 ALFWorld 和 WebShop 任务中，仅需单次或两次 ReAct 提示，模型的性能就能超越经过 $10^3{\sim}10^5$ 个任务实例训练的模仿学习或强化学习方法，成功率分别实现了 34\% 和 10\% 的绝对提升。此外，通过对比纯行动的受控基线模型，我们验证了稀疏、灵活的推理在决策中的重要性，ReAct 在这一方面展现出了持续的优势。除了通用性强、性能提升显著外，推理与行动的融合还让模型在所有任务领域中，都具备更好的可解释性、可信度和可诊断性：人类能够清晰区分模型的结论是来自内部知识还是外部环境，也能通过检查推理轨迹，理解模型做出行动决策的依据。

本研究的核心贡献如下：（1）提出 ReAct 这一全新的基于提示的范式，实现大语言模型中推理与行动的协同，适用于通用任务求解；（2）在多种基准数据集上开展大量实验，验证了在少样本学习设置下，ReAct 相比单独进行推理或行动生成的传统方法具有显著优势；（3）通过系统性的消融实验与分析，阐明了行动在推理任务中的作用，以及推理在交互式任务中的价值；（4）分析了基于提示设置的 ReAct 存在的局限性（即对推理和行动行为的支撑能力有限），并通过初步的微调实验，验证了 ReAct 在增加训练数据后具备进一步性能提升的潜力。将 ReAct 扩展至更多任务的训练与部署，并与强化学习等互补范式相结合，有望进一步释放大语言模型的潜能。

\section{ReAct：推理与行动的协同机制}

考虑智能体与环境交互以完成任务的通用场景：在时间步 $t$，智能体从环境中接收观测信息 $o_t \in \mathcal{O}$，并根据策略 $\pi(a_t \mid c_t)$ 执行行动 $a_t \in \mathcal{A}$，其中 $c_t = (o_1, a_1, \cdots, o_{t-1}, a_{t-1}, o_t)$ 是智能体的上下文信息。当上下文到行动的映射 $c_t \mapsto a_t$ 高度隐含且需要大量计算时，学习该策略将极具挑战性。

ReAct 的核心思想十分简洁：我们将智能体的行动空间扩展为 $\hat{\mathcal{A}} = \mathcal{A} \cup \mathcal{L}$，其中 $\mathcal{L}$ 为语言空间。语言空间中的行动 $\hat{a}_t \in \mathcal{L}$（我们将其称为思维或推理轨迹）不会对外部环境产生影响，因此也不会得到观测反馈。相反，思维 $\hat{a}_t$ 的作用是通过对当前上下文 $c_t$ 进行推理，整合有用信息，并更新上下文 $c_{t+1} = (c_t, \hat{a}_t)$，为后续的推理或行动提供支撑。有效的思维形式具有多样性，例如：拆解任务目标并制定行动计划、融入与任务求解相关的常识知识、从观测信息中提取关键内容、追踪任务进度并调整行动计划、处理突发状况并修正行动计划等。

但由于语言空间 $\mathcal{L}$ 具有无限性，在这种扩展的行动空间中进行学习难度较大，且需要强大的语言先验知识。本研究主要聚焦于冻结大语言模型的设置：以 PaLM-540B 为基础模型，通过少样本的上下文示例进行提示，让模型同时生成特定领域的行动和自由形式的语言思维，以完成任务。每个上下文示例都是人类为解决某一任务实例所生成的行动、思维与环境观测的轨迹。

由于决策与推理能力被整合进同一个大语言模型，ReAct 具备以下四大独特特征：

\textbf{A）设计直观、易于实现：} ReAct 的提示词设计十分简便，标注人员仅需在自身的行动步骤旁，用语言记录下推理思路即可。本研究未使用任何特殊的格式设计、思维模板或示例筛选策略。

\textbf{B）通用性强、灵活适配：} 得益于灵活的思维空间和思维-行动生成格式，ReAct 可适用于各类具有不同行动空间和推理需求的任务，包括但不限于问答、事实验证、文本游戏和网页导航。

\textbf{C）性能优异、鲁棒性高：} 仅通过 1--6 个上下文示例进行学习，ReAct 就能对新的任务实例展现出良好的泛化能力，在不同领域中，性能持续优于仅进行推理或仅进行行动的基线模型。

\textbf{D）人机对齐、可调控性：} ReAct 实现了可解释的序列决策与推理过程，人类能够轻松检查推理过程的合理性和事实的正确性。此外，人类还能通过编辑思维内容，实时控制或修正智能体的行为。

\section{知识密集型推理任务}

我们首先将 ReAct 应用于多跳问答、事实验证等知识密集型推理任务。通过与维基百科 API 交互，ReAct 能够检索信息为推理提供支撑，同时通过推理确定下一步的检索目标，实现了推理与行动的协同。

\subsection{实验设置}

\subsubsection*{任务领域}

我们选取两个对知识检索和推理能力要求较高的数据集展开实验：（1）HotPotQA，一个多跳问答基准数据集，要求模型对至少两段维基百科文本进行推理以得出答案；（2）FEVER，一个事实验证基准数据集，根据是否存在可验证该声明的维基百科文本，将每个声明标注为"支持""反驳"或"信息不足"。本研究中，两个任务均采用仅问题输入的设置：模型仅接收问题/声明作为输入，无法获取支撑文本，必须依靠内部知识，或通过与外部环境交互检索知识来完成推理。

\subsubsection*{行动空间}

我们设计了一个简易的维基百科网页 API，包含三类行动，以支持交互式信息检索：（1）\texttt{search[实体]}：若该实体的维基百科页面存在，返回页面的前 5 个句子；若不存在，由维基百科搜索引擎推荐相似度最高的 5 个实体。（2）\texttt{lookup[字符串]}：返回页面中包含该字符串的下一个句子，模拟浏览器的"查找（Ctrl+F）"功能。（3）\texttt{finish[答案]}：提交答案并完成当前任务。

\subsection{实验方法}

\subsubsection*{ReAct 提示方法}

在 HotPotQA 和 FEVER 任务中，我们分别从训练集中随机选取 6 个和 3 个任务实例，人工构建 ReAct 格式的轨迹，作为提示词中的少样本示例。每个轨迹由多个"思维-行动-观测"步骤构成（即密集型思维），自由形式的思维被用于多种推理场景，具体包括：拆解问题、从维基百科观测信息中提取关键内容、进行常识推理或算术推理、指导检索策略调整，以及整合信息得出最终答案。

\subsubsection*{基线模型}

我们通过对 ReAct 轨迹进行系统性的消融，构建了多种基线模型的提示词：（a）标准提示（Standard）：移除 ReAct 轨迹中的所有思维、行动和观测信息。（b）思维链提示（CoT）：移除 ReAct 轨迹中的行动和观测信息，作为仅推理的基线模型。我们还构建了自洽性思维链基线（CoT-SC）：推理阶段将解码温度设为 0.7，采样 21 条思维链轨迹，选取出现次数最多的答案作为最终结果。（c）仅行动提示（Act）：移除 ReAct 轨迹中的思维信息。

\subsubsection*{内部与外部知识的融合策略}

我们观察到，ReAct 的问题求解过程更贴合事实、更具实际依据，而 CoT 在构建推理结构方面更准确，但容易产生事实幻觉或不合理的推理。因此，我们提出将 ReAct 与 CoT-SC 相结合的策略，让模型根据以下启发式规则，自主决定切换至另一种方法：A）ReAct $\to$ CoT-SC：若 ReAct 在设定步数内未得出答案，则切换至 CoT-SC。B）CoT-SC $\to$ ReAct：若 $n$ 条 CoT-SC 采样轨迹中，出现次数最多的答案的占比不足 $n/2$（即模型的内部知识无法为任务提供足够可靠的支撑），则切换至 ReAct。

\subsection{实验结果与分析}

\subsubsection*{ReAct 性能持续优于仅行动模型}

ReAct 在两个任务上的性能均优于仅行动模型，证明了推理对行动的指导价值 —— 尤其是在整合信息得出最终答案的阶段。微调实验的结果也进一步证实，推理轨迹能让模型的行动决策更具依据。

\subsubsection*{ReAct 与 CoT 的性能对比}

ReAct 在 FEVER 任务上的性能优于 CoT（60.9 vs.\ 56.3），在 HotPotQA 任务上的性能略低于 CoT（27.4 vs.\ 29.4）。主要发现如下：

A）事实幻觉是 CoT 的核心问题：在成功求解的案例中，CoT 的假阳性率远高于 ReAct（14\% vs.\ 6\%）；而事实幻觉也是 CoT 最主要的失败原因，占比达 56\%。相比之下，由于能访问外部知识库，ReAct 的问题求解过程更贴合实际、以事实为驱动，可信度更高。

B）思维-行动-观测的交替结构存在局限性：这种结构虽提升了 ReAct 的事实依据性和可信度，但也降低了其构建推理步骤的灵活性，导致 ReAct 的推理错误率高于 CoT。ReAct 存在一种典型的错误模式：模型会重复生成之前的思维和行动，无法推理出合理的下一步行动，也无法跳出循环。

C）有效检索是 ReAct 成功的关键：无信息价值的检索是 ReAct 的主要失败原因之一，占比 23\%。这一现象可视为模型事实性与灵活性的权衡。

\subsubsection*{ReAct 与 CoT-SC 融合是提示学习的最优策略}

HotPotQA 和 FEVER 任务上性能最优的提示方法分别为 ReAct $\to$ CoT-SC 和 CoT-SC $\to$ ReAct。两种融合策略在所有采样数下均显著且持续优于纯 CoT-SC；仅需 3--5 个采样数，融合策略就能达到纯 CoT-SC 21 个采样数的性能。这些结果表明，将模型的内部知识与外部知识合理融合，对推理任务具有重要价值。

\subsubsection*{ReAct 是微调学习的最优方法}

以 PaLM-8B/62B 为基础模型时，由于从上下文示例中同时学习推理和行动的难度较大，ReAct 在提示学习中的性能排名最后。但仅通过 3000 个示例进行微调后，ReAct 成为四种方法中性能最优的模型：微调后的 PaLM-8B-ReAct 性能优于所有 PaLM-62B 的提示学习模型，微调后的 PaLM-62B-ReAct 性能甚至优于所有 PaLM-540B 的提示学习模型。这表明，学习如何通过推理行动获取信息，是一种更具泛化性的知识推理能力。

\section{决策任务}

我们还将 ReAct 应用于两类基于语言的交互式决策任务 —— ALFWorld 和 WebShop。这两个任务的环境均具有复杂性，要求智能体在稀疏奖励的条件下，完成长序列的行动决策，因此亟需通过推理为行动和探索提供有效指导。

\subsubsection*{ALFWorld 任务}

ALFWorld 是一个合成的文本基游戏环境，与具身智能体基准数据集 ALFRED 相适配。该环境包含 6 类任务，要求智能体通过文本行动与模拟家居环境交互、导航，完成高层目标（如"在台灯下查看纸张"）。每个任务实例包含超过 50 个场景位置，专家策略需要执行 50 余步才能完成，因此对智能体的子目标规划、追踪能力，以及系统性探索能力提出了很高要求。ALFWorld 的一个核心挑战，是要求智能体判断家居物品的可能位置（如"台灯通常在桌子、架子或梳妆台上"），这一特点让大语言模型能够充分利用预训练的常识知识。

在 ReAct 提示词设计中，我们为每类任务从训练集中随机标注 3 条轨迹，每条轨迹包含稀疏的思维，主要用于：（1）拆解总目标；（2）追踪子目标的完成情况；（3）确定下一个子目标；（4）通过常识推理判断物品的可能位置及操作方式。

\subsubsection*{WebShop 任务}

WebShop 是一个近期提出的在线购物网页环境，包含 118 万件真实商品和 1.2 万条人类指令。该任务要求智能体通过网页交互，根据用户指令选购商品（如"我想找一个带抽屉的床头柜，表面为镍质工艺，价格低于 140 美元"）。该任务的评估指标为平均得分（所有任务中，选中商品满足用户需求属性的比例均值）和成功率。

\subsubsection*{实验结果}

ReAct 在 ALFWorld 和 WebShop 两个任务上的性能均优于仅行动模型。在 ALFWorld 任务中，ReAct 的最优实验结果实现了 71\% 的平均成功率，显著优于仅行动模型的最优结果（45\%）和 BUTLER 模型（37\%）；即便 ReAct 的最差实验结果（48\%），也优于后两者的最优结果。在 WebShop 任务中，增加了稀疏推理的 ReAct，成功率相比此前的最优结果实现了 10\% 的绝对提升。

\subsubsection*{内部推理与外部反馈的价值}

ReAct 是首个将大语言模型的推理与行动融合，应用于闭环系统交互式环境的研究。与之最接近的前期研究是《内心独白》，但 ReAct 在决策任务中的推理轨迹更灵活、更稀疏，能够针对不同任务，生成多种类型的推理。消融实验证明 ReAct 的整体成功率（71\%）显著优于 IM 式提示模型（53\%）。

\section{相关工作}

\subsubsection*{大语言模型的推理研究}

利用大语言模型进行推理的最经典研究，是思维链（CoT） —— 该研究首次发现，大语言模型能够生成推理轨迹，解决算术、常识等推理问题。后续相关研究层出不穷，包括用于解决复杂任务的由浅入深提示法、零样本思维链，以及自洽性推理。近期研究系统性地研究了思维链的构成与结构，发现符号、模式和文本的结合，是思维链有效的关键。

还有研究突破了简单的提示学习，提出了更复杂的推理架构：例如选择-推理（Selection-Inference）将推理过程拆解为"选择"和"推理"两个步骤；自举推理（STaR）通过在模型自身生成的正确推理过程上微调，实现推理能力的自举提升；可信推理（Faithful Reasoning）将多步推理拆解为三个步骤，每个步骤由专用的大语言模型完成。

与上述方法不同，ReAct 的核心并非孤立、固定的推理，而是将模型的行动及对应的观测信息，融入连贯的输入流中，让模型做出更准确的推理，并解决推理之外的任务（如交互式决策）。

\subsubsection*{大语言模型的决策研究}

大语言模型的强大能力，使其能够完成语言生成之外的任务；将大语言模型作为策略模型应用于决策任务，尤其是交互式环境中的决策，正成为研究热点。WebGPT 利用大语言模型与浏览器交互、导航网页，并从 ELI5 数据集中推导复杂问题的答案。与 ReAct 不同，WebGPT 并未显式建模思考与推理过程，而是依靠高成本的人类反馈进行强化学习。

大语言模型也越来越多地被应用于交互式和具身环境的规划与决策，其中与 ReAct 最相关的研究是 SayCan 和《内心独白》，二者均将大语言模型用于机器人的行动规划与决策。与 ReAct 不同，这两项工作均未真正实现自由形式的内部推理，其"内心独白"本质上是对环境状态的外部反馈注入。

\section{结论}

本研究提出了 ReAct —— 一种简洁且有效的方法，实现了大语言模型中推理与行动的协同。通过在多跳问答、事实验证和交互式决策任务上开展大量实验，我们验证了 ReAct 能够在生成可解释决策轨迹的同时，实现性能的提升。尽管 ReAct 方法简洁，但面对具有大行动空间的复杂任务，模型需要更多的示例才能充分学习；而这一需求，很容易超出上下文学习的输入长度限制。我们在 HotPotQA 任务上探索了微调方法，取得了初步的积极结果；而利用更多高质量的人工标注数据进行学习，将是进一步提升 ReAct 性能的关键。将 ReAct 扩展至多任务训练，并与强化学习等互补范式相结合，有望构建性能更优的智能体，进一步释放大语言模型在更多实际应用中的潜能。

\begingroup
    \linespreadsingle{}
    \printbibliography[title={外文翻译参考文献}]
\endgroup
